\subsection{Recommendation Performance}
Table~\ref{tab:main_results} reports HR@10 and NDCG@10 (99 negatives, seed 42).

\paragraph{Pipeline~A (Behavioral Graph + BGE-M3).}
HR@10 = \textbf{0.9047} (NDCG@10 = 0.5393), a $+108.2\%$ improvement over
Content-KNN (0.4346).

\paragraph{DIF-SASRec (Pipeline~B).}
HR@10 = 0.7745 (NDCG@10 = 0.5024), $+78.2\%$ over Content-KNN.

\paragraph{Dual-pipeline system.}
The union achieves \textbf{HR@10 = 0.9736} (NDCG@10 = 0.5571).
Pipeline~A rescues 88.3\% of DIF-SASRec misses; DIF-SASRec rescues 72.3\%
of Pipeline~A misses. Pipeline~A excels on items within the behavioral
graph index (375k items), while DIF-SASRec covers items outside it.

\input{tables/main_results}
\input{tables/complementarity}

\subsection{Robustness Under Stricter Evaluation}
\label{ssec:robustness}

Table~\ref{tab:robustness} repeats the comparison under 999 negatives
(random-chance HR@10 = 0.01).

Pipeline~A achieves HR@10 = 0.3272; DIF-SASRec leads on NDCG@10 (0.2209 vs.\
0.2102). Content-KNN drops to 0.2122. The relative advantage over
chance grows with negative pool size for both pipelines.

\begin{table}[tbp]
\caption{Robustness evaluation on Amazon Books~\cite{hou2024amazon} (100k users, 999 random
         negatives per user; random-chance baseline HR@10 = 0.01).
         Best results in \textbf{bold}.
         The union system (A $\cup$ B) is excluded from this table; it is
         evaluated under the 99-negative protocol in Table~\ref{tab:main_results}.
         This table isolates individual pipeline robustness.}
\label{tab:robustness}
\centering
\begin{tabular}{lcc}
\toprule
\textbf{Model} & \textbf{HR@10} & \textbf{NDCG@10} \\
\midrule
Content-KNN (BGE-M3)       & 0.2122          & 0.1349          \\
DIF-SASRec (Pipeline~B)    & 0.3142          & \textbf{0.2209} \\
Pipeline~A (Cleora+BGE-M3) & \textbf{0.3272} & 0.2102          \\
\bottomrule
\end{tabular}
\end{table}


\subsection{Ablation Study}
Table~\ref{tab:ablation} reports both pipeline-level and component-level ablations.
Pipeline-component scores are scaled proportionally to enable direct comparison
with pipeline-level rows, preserving relative $\Delta$HR.

\paragraph{Pipeline ablation.}
Pipeline~A alone (+16.8\% HR@10 over Pipeline~B, +108.2\% over Content-KNN)
shows that behavioral graph candidates with BGE-M3 profile ranking produce
a strong signal. The union adds +7.6\% by covering items outside the graph.

\paragraph{Component ablation.}
Removing the category head reduces HR@10 by $1.8\%$; disabling the content
veto ($\tau=0.3$) in Pipeline~B costs $2.1\%$; replacing decoupled attention
with a single stream costs $4.6\%$; and removing per-user fine-tuning causes
the largest drop ($-7.9\%$), confirming online adaptation as the primary
personalization driver.

\begin{table}[tbp]
\caption{Ablation study on the Amazon Books~\cite{hou2024amazon} test split (100k users,
         99-negative sampled evaluation, seed 42).
         \emph{Top}: pipeline-level ablation, standalone and union results.
         \emph{Bottom}: DIF-SASRec component ablation; $\Delta$HR is relative
         to the DIF-SASRec full model (see Section~\ref{sec:exp} for scaling note).}
\label{tab:ablation}
\centering
\begin{tabular}{lccc}
\toprule
\textbf{Configuration} & \textbf{HR@10} & \textbf{NDCG@10} & \textbf{$\Delta$HR} \\
\midrule
\multicolumn{4}{l}{\textit{Pipeline ablation}} \\
Pipeline~B only (DIF-SASRec)  & 0.7745 & 0.5024 & --- \\
Pipeline~A only (Cleora)      & 0.9047 & 0.5393 & $+16.8\%$ \\
System (A $\cup$ B)           & \textbf{0.9736} & \textbf{0.5571} & $+25.7\%$ \\
\midrule
\multicolumn{4}{l}{\textit{DIF-SASRec component ablation}} \\
DIF-SASRec (full)              & 0.7745 & 0.5024 & --- \\
w/o category auxiliary loss   & 0.7608 & 0.4889 & $-1.8\%$ \\
w/o content veto (Pipeline~B) & 0.7584 & 0.4814 & $-2.1\%$ \\
w/o online adaptation         & 0.7129 & 0.4467 & $-7.9\%$ \\
w/o dual-stream attention     & 0.7389 & 0.4683 & $-4.6\%$ \\
\bottomrule
\end{tabular}
\end{table}


\subsection{Latency Analysis}
Table~\ref{tab:latency} reports end-to-end latency under two conditions:
(i)~\emph{cold}, 35 unique queries each run once (guaranteed cache misses);
(ii)~\emph{warm}, repeated queries after the HNSW index is OS-page-cached
and LRU caches are populated.

\input{tables/latency}

\paragraph{Cold path.}
Cold P50 is 1\,127\,ms (English) and 1\,281\,ms (non-English).
The bottleneck is FAISS HNSW (P50 906\,ms) with OS page faults on the
12.6\,GB index. NLLB translation adds 156\,ms for non-English.

\paragraph{Warm path.}
Once warm, P50 drops to 17\,ms (internal) / 34\,ms (E2E) for English.
FAISS is 2\,ms; translation and encoding are served from LRU caches
at under 1\,ms. BM25 (${\approx}6$\,ms) is the bottleneck.

\paragraph{Production characterization.}
A startup warmup pages the HNSW index into RAM. New unique queries
incur encoding (${\approx}40$\,ms) plus FAISS (${\approx}2$\,ms),
steady-state P50 60--80\,ms. Cached queries serve at ${<}\,25$\,ms.

\subsection{Concurrency and Memory}
The pool of 8 DIF-SASRec instances sustains ${\geq}10$ concurrent users
(total footprint 1.18\,GB). Each instance holds pretrained weights in CPU
RAM; fine-tuned weights occupy GPU memory only per request.

\subsection{Encoder Comparison: BLaIR vs.\ BGE-M3}
\label{ssec:encoder}

We evaluate the encoder choice that underpins the active-search and
passive-recommendation pipelines.
Table~\ref{tab:encoder_comparison} compares the legacy BLaIR
encoder~\cite{hou2024amazon} against BGE-M3~\cite{chen2024bgem3} along two axes.
\input{tables/encoder_comparison}

\paragraph{Sampled evaluation.}
Under 99-negative evaluation, BLaIR shows higher HR@10 (0.5484 vs.\ 0.4510),
but the benchmark is confounded: clicks were collected while BLaIR was
the live encoder, so held-out items correlate with its embedding space.
Score analysis confirms this: BLaIR's intra- and inter-genre cosine
similarities (0.819 vs.\ 0.765, ratio 1.061) are nearly indistinguishable,
indicating embedding collapse~\cite{gao2021simcse}. BGE-M3's higher ratio
(1.147) reflects distinct genre clusters.

\paragraph{Query retrieval quality.}
Live evaluation with 22 labeled queries shows BGE-M3 leading on all
groups except single-token English: medium (0.333 vs.\ 0.083), long
(0.400 vs.\ 0.200), non-English (0.275--0.550 vs.\ 0.025--0.050).
BLaIR leads on short queries (0.233 vs.\ 0.150), likely due to collapsed
similarities returning genre-coherent results regardless of the token.
We select BGE-M3 as the production encoder given typed text search
is the primary interaction mode and BLaIR's sampled-evaluation advantage
is confounded.
